\begin{table}[htp]
\centering

    % Esta primera línea define las columnas de la tabla. Los posibles tipos de columna son:
    % c: texto centrado
    % l: texto alineado a la izquierda
    % r: texto centrado a la derecha
    % p: columna de ancho fijo
    % Las columnas tienen ancho dinámico según la anchura máxima de los elementos que contengan.

    % Las columnas l/r/c no parten el texto en filas diferentes si éste es demasiado largo. Para ello, puede utilizar el tipo de columna de ancho fijo "p".
    
    % Las barras verticales | se usan para definir los bordes verticales de la tabla. Pruebe a eliminar algunas y observe qué ocurre.
    \begin{center}
    \begin{tabular}{ | l | l | }
        
        % A continuación van las filas de la tabla. En cada fila, las columnas se separan con el carácter &
        % Para terminar una fila se usa \\
        % Para incluir un borde horizontal entre filas se usa \hline
        
        % Cabecera con textos en negrita:
        \hline
        \textbf{Media}  & 0.132385 \\
        \hline
        \textbf{Mediana} & 0.130000 \\
        \hline
        \textbf{Moda} & 0.125230 \\
        \hline
        \textbf{IQR} & 0.010000 \\
        \hline
        \textbf{Varianza} & 0.001912 \\
        \hline
        \textbf{Desviación estándar} & 0.043723 \\
        \hline
        \textbf{Covarianza (v)} & 3.647974 \\
        \hline
        \textbf{Correlación (v)} & 0.289460 \\
        \hline
        \textbf{Covarianza (GPU)} & 0.001354 \\
        \hline
        \textbf{Correlación (GPU)} & 0.759181 \\
        \hline
        \textbf{Q1} & 0.13 \\
        \hline
        \textbf{Q2} & 0.13 \\
        \hline
        \textbf{Q3} & 0.14 \\
        \hline
        
    \end{tabular} 
    \end{center}
    
    \caption{Resultados de la renderización en la prueba del polígono: \textbf{CPU}}
    
    \label{tbl:results-polygon-render-cpu}
\end{table}